\chapter{Solutions: Midterm, 5 December 2014}
\label{s:theory-2014-12-midterm}
\begin{sourceNote}
All seven printed questions survive. Their weights total 110\% and are retained. The derivations below are independent of handwritten answers. In Question 2, $a$ is the conventional cubic lattice parameter; in Question 3 it is the two-dimensional primitive-translation length.
\end{sourceNote}

\section*{1. Eight definitions}
\begin{enumerate}[label=(\alph*)]
\item \textbf{Bravais lattice.}
A three-dimensional Bravais lattice is the point set
\[
\mathcal L=\{\vct R=n_1\vct t_1+n_2\vct t_2+n_3\vct t_3:n_i\in\mathbb Z\},
\]
with linearly independent primitive vectors. Every lattice point has the same environment under translation. A lattice describes translational symmetry; a lattice point need not represent one atom.
\item \textbf{Composite crystal.}
In the lattice-plus-basis terminology relevant here, attach the same atomic basis $\{(\vct d_\alpha,s_\alpha)\}$ to each Bravais point. Positions are $\vct R+\vct d_\alpha$, with $s_\alpha$ specifying species. Even identical atoms at different basis sites need not be related by a Bravais translation. The honeycomb crystal has a two-atom basis on a triangular Bravais lattice.
\item \textbf{Brillouin zone.}
The first Brillouin zone is the reciprocal origin's Wigner--Seitz cell:
\[
|\vct k|\leq|\vct k-\vct G|\quad\hbox{for every }\vct G\in\mathcal L^*.
\]
Its faces are perpendicular bisectors between the origin and other reciprocal points. Higher zones can be classified by the number of crossed Bragg planes.
\item \textbf{Reciprocal lattice.}
With $\vct t_i\cdot\vct g_j=2\pi\delta_{ij}$,
\[
\mathcal L^*=\{m_1\vct g_1+m_2\vct g_2+m_3\vct g_3:m_i\in\mathbb Z\}.
\]
Equivalently $\ee^{\ii\vct G\cdot\vct R}=1$ for every direct-lattice translation.
\item \textbf{Bloch theorem.}
For a single-particle Hamiltonian with lattice symmetry, eigenstates can be chosen as
\[
\psi_{n\vct k}(\vct r)=\ee^{\ii\vct k\cdot\vct r}u_{n\vct k}(\vct r),\quad
u_{n\vct k}(\vct r+\vct R)=u_{n\vct k}(\vct r).
\]
Thus $\psi(\vct r+\vct R)=\ee^{\ii\vct k\cdot\vct R}\psi(\vct r)$: $\psi$ is generally quasiperiodic, while $u$ is periodic. The theorem follows from simultaneous diagonalization of the Hamiltonian and commuting translations.
\item \textbf{Exchange energy.}
For a determinant of orthonormal spin orbitals $\chi_i(x)$, $x=(\vct r,\sigma)$,
\[
E_{ee}=\frac12\sum_{ij}(J_{ij}-K_{ij}),\qquad E_x=-\frac12\sum_{ij}K_{ij},
\]
\[
K_{ij}=\iint\chi_i^*(x)\chi_j^*(x')
\frac{\kappa\,\chi_j(x)\chi_i(x')}{|\vct r-\vct r'|}\dd x\dd x',
\qquad\kappa=\frac{e^2}{4\pi\epsilon_0}.
\]
Exchange arises from antisymmetry; separated spin functions restrict it to equal-spin pairs. It is not an additional classical force or the entire correlation energy.
\item \textbf{Fully spin-polarized ground state.}
All spins point in the same direction: $N_\uparrow=N$, $N_\downarrow=0$ gives $\zeta=(N_\uparrow-N_\downarrow)/N=1$, $S_z=N\hbar/2$, and total-spin quantum number $S=N/2$. It does not mean zero spin.
\item \textbf{Wigner crystallization.}
At sufficiently low density, an electron gas in a uniform positive background can form an ordered electron crystal. With mean spacing $r_sa_0$, kinetic and Coulomb scales vary as $r_s^{-2}$ and $r_s^{-1}$, respectively; interactions become relatively more important at large $r_s$. This is spontaneous density ordering, not merely drawing a Wigner--Seitz cell.
\end{enumerate}

\section*{2. Lattice planes and translation phase}
\subsection*{Reciprocal basis and primitive volume}
From $\vct g_1=2\pi(\vct t_2\times\vct t_3)/\Omega$ and cyclic permutations,
\[
\Omega=\vct t_1\cdot(\vct t_2\times\vct t_3)=\frac{a^3}{4},
\]
\[
\boxed{\vct g_1=\frac{2\pi}{a}(-1,1,1),\quad
\vct g_2=\frac{2\pi}{a}(1,-1,1),\quad
\vct g_3=\frac{2\pi}{a}(1,1,-1).}
\]
One point per primitive cell gives $\rho_3=4/a^3$. For coprime integers $h_i$, let $\vct G_0=\sum_i h_i\vct g_i$. At lattice points,
\[
\vct G_0\cdot\vct R=2\pi(h_1n_1+h_2n_2+h_3n_3)=2\pi m.
\]
Every integer $m$ occurs. Thus successive planes containing lattice points have
\[
\boxed{d=2\pi/|\vct G_0|.}
\]
Use the shortest reciprocal vector parallel to the specified normal; otherwise the level planes may include empty planes.

\subsection*{(a) Normal $3\vct g_1+6\vct g_2-6\vct g_3$}
Remove the common factor $3$:
\[
\vct G_0=\vct g_1+2\vct g_2-2\vct g_3
=\frac{2\pi}{a}(-1,-3,5),\qquad\gcd(1,2,-2)=1.
\]
Consequently
\[
\boxed{d=a/\sqrt{35}.}
\]
Directly using the printed normal's length would incorrectly divide this answer by three.

\subsection*{(b) Surface density}
A primitive planar cell of area $A_2$, stacked by a lattice translation with normal component $d$, gives $A_2d=\Omega$. Hence
\[
\boxed{\rho_2=\frac1{A_2}=\frac d\Omega=\frac4{a^2\sqrt{35}}.}
\]
For an explicit check, all integer solutions of $n_1+2n_2-2n_3=0$ are generated by
\[
\vct u=-2\vct t_1+\vct t_2=\frac a2(1,-2,-1),\qquad
\vct v=2\vct t_1+\vct t_3=\frac a2(1,3,2).
\]
Their cross-product magnitude is $A_2=a^2\sqrt{35}/4$.

\subsection*{(c) Normal $\tfrac12\vct t_1+\vct t_2-\vct t_3$}
The direction is $a(0,-1,3)/4$. A reciprocal vector has components
\[
\frac{2\pi}{a}(-h_1+h_2+h_3,\ h_1-h_2+h_3,\ h_1+h_2-h_3).
\]
Setting these equal to $(2\pi/a)(0,-\lambda,3\lambda)$ gives $(h_1,h_2,h_3)=(\lambda,3\lambda/2,-\lambda/2)$. The smallest positive $\lambda$ giving integer coefficients is $2$:
\[
\vct G_0=2\vct g_1+3\vct g_2-\vct g_3=\frac{2\pi}{a}(0,-2,6),
\qquad\boxed{d=a/\sqrt{40}.}
\]
The shorter vector $(2\pi/a)(0,-1,3)$ is not reciprocal: its primitive coefficients are half-integers.

\subsection*{(d) Periodicity}
For $f(\vct r)=\exp(\ii\vct g_1\cdot\vct r/2)$,
\[
f(\vct r+\vct t_1)=f(\vct r)\ee^{\ii\pi}=-f(\vct r).
\]
Thus \textbf{no}. The squared magnitude is uniform, but the complex wavefunction retains only translations with even $n_1$, not all periods of the original lattice.

\section*{3. Critical points and two-dimensional DOS}
Let $E_\pm=E_p\pm\gamma\sqrt S$, with
\[
S=1+4\cos x\cos y+4\cos^2x,\quad x=\frac{ak_x}{2},\quad y=\frac{\sqrt3ak_y}{2}.
\]
An ordinary stationary critical point is differentiable and satisfies $\nabla E=0$. At a touching, $\nabla S=0$ does not by itself imply $\nabla\sqrt S=0$.

\subsection*{(a) $\Gamma$ and $P$}
At $\Gamma$,
\[
S=9-6x^2-2y^2+O(k^4)=9-\frac32a^2k^2+O(k^4),
\]
\[
\sqrt S=3-\frac{a^2k^2}{4}+O(k^4),\qquad
E_\pm=E_p\pm3\gamma\mp\frac{\gamma a^2}{4}k^2+O(k^4).
\]
Hence $\Gamma$ is a maximum of the upper band and a minimum of the lower band. The Hessians $\mp(\gamma a^2/2)I$ are definite.

At $P=(4\pi/(3a),0)$, write $\vct k=\vct P+\vct q$. Then
\[
S=\frac34a^2(q_x^2+q_y^2)+O(q^3),\qquad
\boxed{E_\pm=E_p\pm\beta|\vct q|+O(q^2),\quad\beta=\frac{\sqrt3\gamma a}{2}.}
\]
Along $q_y=0$, the term $\beta|q_x|$ has unequal left and right derivatives. Therefore $P$ is a \textbf{Dirac touching, not a differentiable zero-gradient critical point}. A broader terminology may call it a nonsmooth upper-band minimum and lower-band maximum, but the conical nature must be stated; a quadratic effective-mass approximation is invalid.

\begin{center}
\begin{tikzpicture}[x=1cm,y=0.75cm,>=Stealth]
\draw[->] (-2.3,0)--(2.4,0) node[right] {$q_x$};
\draw[->] (0,-2.4)--(0,2.5) node[above] {$E-E_p$};
\draw[very thick,blue] (-2,2)--(0,0)--(2,2);
\draw[very thick,red!70!black] (-2,-2)--(0,0)--(2,-2);
\node[blue,right] at (1.55,1.8) {$E_+$};
\node[red!70!black,right] at (1.55,-1.8) {$E_-$};
\node[below left] at (0,0) {$P$};
\end{tikzpicture}

\small Section through a Dirac cone: opposite sides have different slopes.
\end{center}

\subsection*{(b) DOS around the lower-band minimum at $\Gamma$}
Define the DOS per real-space area, per spin, for one local patch:
\[
D(E)=\int_{\rm patch}\frac{\dd^2q}{(2\pi)^2}\delta(E-E(\vct q)).
\]
Let $E_0=E_p-3\gamma$, $\alpha=\gamma a^2/4$, and use $E=E_0+\alpha q^2$ in $q<\Lambda$. Then
\begin{align*}
D_-(E)&=\frac1{2\pi}\int_0^\Lambda q\,\dd q\,\delta(E-E_0-\alpha q^2)\\
&=\boxed{\frac1{4\pi\alpha}\Theta(E-E_0)\Theta(E_0+\alpha\Lambda^2-E)}\\
&=\frac1{\pi\gamma a^2}\Theta(E-E_0)\Theta(E_0+\alpha\Lambda^2-E).
\end{align*}
This is the constant step at a two-dimensional parabolic extremum. The upper-band maximum $E_1=E_p+3\gamma$ gives a reversed step with the same height for $E_1-\alpha\Lambda^2<E<E_1$. Multiply by sample area for total DOS and by two for both spins. No extra valley factor belongs to a single patch.

For comparison, one Dirac valley per spin gives
\[
D_P(E)=\frac{|E-E_p|}{2\pi\beta^2}=\frac{2|E-E_p|}{3\pi\gamma^2a^2},
\]
which vanishes at the touching energy.

\section*{4. Hamiltonian with fixed nuclei}
Fix nuclear coordinates $\vct R_A$ and charges $Z_Ae$. Neglect relativity, spin--orbit coupling, and nuclear motion. Define
\[
v_{\rm ext}(\vct r)=-\sum_A\frac{Z_A\kappa}{|\vct r-\vct R_A|},\qquad
h(i)=-\frac{\hbar^2}{2m_e}\nabla_i^2+v_{\rm ext}(\vct r_i).
\]
Then
\[
\boxed{H_e=\sum_{i=1}^Nh(i)+\sum_{i<j}\frac{\kappa}{r_{ij}}.}
\]
These terms are electron kinetic energy, nuclear attraction, and electron repulsion. Total energy at the fixed nuclear configuration also includes the constant
\[
E_{nn}=\sum_{A<B}\frac{Z_AZ_B\kappa}{|\vct R_A-\vct R_B|}.
\]
This does not affect electronic variation. Writing the pair sum as $\sum_{i\ne j}$ requires a factor $1/2$.

\section*{5. Hartree trial state}
For $x_i=(\vct r_i,\sigma_i)$,
\[
\boxed{\Psi_H(x_1,\ldots,x_N)=\prod_{i=1}^N\chi_i(x_i),\qquad
\int|\chi_i(x)|^2\dd x=1.}
\]
The expectation value is
\[
E_H=\sum_i h_{ii}+\frac12\sum_{i\ne j}J_{ij},\qquad
J_{ij}=\iint\frac{\kappa|\chi_i(x)|^2|\chi_j(x')|^2}{|\vct r-\vct r'|}\dd x\dd x'.
\]
Variation of each normalized orbital gives the average Hartree field of the other electrons. The product does not enforce antisymmetry, so it does not fully satisfy fermionic statistics and has no exchange term.

\section*{6. Hartree--Fock trial state and optimization}
\subsection*{Antisymmetrization and energy}
Take orthonormal spin orbitals:
\[
\boxed{\Psi_{\rm HF}(x_1,\ldots,x_N)=\frac1{\sqrt{N!}}\det[\chi_i(x_j)]_{i,j=1}^N,\quad
\braket{\chi_i}{\chi_j}=\delta_{ij}.}
\]
Exchanging two electron coordinates exchanges columns and changes the sign. Its energy is
\[
E_{\rm HF}=\sum_i h_{ii}+\frac12\sum_{ij}(J_{ij}-K_{ij}).
\]
Since $J_{ii}=K_{ii}$, self-interaction cancels.

\subsection*{Constrained variation}
Minimize
\[
\mathcal L=E_{\rm HF}-\sum_{ij}\lambda_{ij}
(\braket{\chi_i}{\chi_j}-\delta_{ij}).
\]
Variation with respect to $\chi_i^*(x)$ gives
\[
F\chi_i=\sum_j\lambda_{ij}\chi_j,\qquad F=h+\sum_{j\in\rm occ}(J_j-K_j),
\]
\begin{align*}
(J_jf)(x)&=\left[\int\frac{\kappa|\chi_j(x')|^2}{|\vct r-\vct r'|}\dd x'\right]f(x),\\
(K_jf)(x)&=\chi_j(x)\int\frac{\kappa\chi_j^*(x')f(x')}{|\vct r-\vct r'|}\dd x'.
\end{align*}
$J_j$ is local; $K_j$ is nonlocal. Diagonalize the Hermitian multiplier matrix by a unitary rotation of occupied orbitals, changing the determinant only by an overall phase. The canonical equations are
\[
\boxed{F\chi_i=\epsilon_i\chi_i.}
\]

\subsection*{Self-consistency and limitations}
Guess orbitals, build the density and Fock operator, solve its eigenproblem, occupy $N$ orbitals, and update until density and energy converge. Compare distinct spin or symmetry solutions; convergence alone does not prove the global minimum. The optimized determinant is a variational approximation in a restricted space, not generally an exact correlated state. Its total energy is
\[
E_{\rm HF}=\sum_i\epsilon_i-\frac12\sum_{ij}(J_{ij}-K_{ij}),
\]
not the uncorrected eigenvalue sum.

\section*{7. Slater--Jastrow trial state}
For a fluid electron gas in periodic volume $\Omega$, use momentum-state determinants times a positive pair-correlation factor:
\[
\boxed{\Psi_T(\vct R)=D_\uparrow D_\downarrow
\exp\left[-\sum_{i<j}u_{\sigma_i\sigma_j}(r_{ij})\right],\qquad
\vct R=(\vct r_1,\ldots,\vct r_N).}
\]
For example $D_\sigma=\det[\ee^{\ii\vct k_\alpha\cdot\vct r_j}/\sqrt\Omega]$ fills the chosen Fermi sea for each spin population. This is a fixed-spin coordinate component; the complete spin wavefunction remains antisymmetric. A fully polarized state has one spin determinant.

The determinants supply antisymmetry and nodes. The symmetric positive Jastrow factor preserves the nodes while changing pair probabilities and lowering correlation energy. Its short-range form can satisfy Coulomb cusp conditions; its long-range form can incorporate screening and collective fluctuations. Optimize its parameters variationally. A trial crystal phase can instead use localized orbitals around lattice sites.

Sampling $p(\vct R)\propto|\Psi_T(\vct R)|^2$ gives the local-energy average:
\[
E_T=\frac{\int|\Psi_T|^2E_L\,\dd\vct R}{\int|\Psi_T|^2\,\dd\vct R},\quad
E_L=\frac{H\Psi_T}{\Psi_T},\quad E_T\geq E_0.
\]
Projection Monte Carlo further filters higher energies with imaginary-time evolution $\ee^{-\tau(H-E_T)}$. Holding trial nodes fixed yields a fixed-node approximation, not an error-free exact result. The Ceperley--Alder stochastic projection study is more than a single variational sampling calculation.

\section*{Primary reference used for checks}
\href{https://journals.aps.org/prl/abstract/10.1103/PhysRevLett.45.566}{D. M. Ceperley and B. J. Alder, Ground State of the Electron Gas by a Stochastic Method, Phys. Rev. Lett. 45, 566 (1980)}.
An \href{https://people.physics.illinois.edu/Ceperley/papers/015.pdf}{author-hosted full text} is available. The lattice geometry, band expansions, DOS factors, and variational steps were derived independently.
